// 323111047

\section{Introducción}

Planteamiento del problema. El problema técnico consiste en desarrollar una aplicación que identifique si una palabra ingresada por el usuario es un palíndromo mediante el uso coordinado de dos estructuras de datos dinámicas. El reto radica en implementar correctamente una pila ligada para recuperar los caracteres en orden inverso y una cola ligada para mantener el orden original de ingreso. La dificultad principal se encuentra en la gestión de nodos y punteros para asegurar que cada carácter de la palabra se almacene simultáneamente en ambas estructuras y que la comparación posterior, elemento por elemento, se realice de forma precisa mediante la extracción coordinada de datos.

\vspace{0.4cm}

Motivación. La importancia de este ejercicio radica en comprender el comportamiento y las diferencias operativas entre las estructuras lineales de datos en un escenario real. Mientras que una cola es fundamental para sistemas de gestión de procesos y flujos constantes, una pila es vital para funciones de reversibilidad y control de estados en memoria. Resolver el problema de los palíndromos mediante estas estructuras permite al estudiante visualizar cómo la organización lógica de la información facilita la resolución de problemas algorítmicos, fortaleciendo el dominio sobre el manejo de memoria dinámica y estructuras ligadas que son esenciales en el desarrollo de software profesional.

\vspace{0.4cm}

Objetivos. El propósito fundamental de esta práctica es lograr la implementación funcional de nodos ligados para construir una pila y una cola de caracteres de manera dinámica. Se busca que el alumno demuestre dominio en las operaciones de inserción y extracción de datos, validando que el flujo de información cumpla con las propiedades de cada estructura. Finalmente, se pretende verificar mediante la comparación de los caracteres obtenidos si la palabra ingresada es un palíndromo, permitiendo así una discusión final sobre la efectividad de las estructuras lineales en la validación de secuencias lógicas.